\documentclass[aps,pra,twocolumn,superscriptaddress,longbibliography,nofootinbib]{revtex4-2}

\usepackage[utf8]{inputenc}
\usepackage[T1]{fontenc}
\usepackage{amsmath,amssymb,amsthm}
\usepackage{bm}
\usepackage{graphicx}
\usepackage{physics}
\usepackage{xcolor}
\usepackage[colorlinks=true,linkcolor=blue,citecolor=blue,urlcolor=blue]{hyperref}

% ---- Convenience macros ----
\newcommand{\Eve}{\mathrm{E}}
\newcommand{\Alice}{\mathrm{A}}
\newcommand{\Bob}{\mathrm{B}}
\newcommand{\pguess}{P_{\Eve}^{\mathrm{guess}}}
\newcommand{\wherekey}{\texttt{where\_key}}
\newcommand{\todo}[1]{\textcolor{red}{[TODO: #1]}}

\begin{document}

\title{Key rates for prepare-and-measure contextuality-based cryptography against operational and quantum entangling-probe adversaries}

\author{Elie Wolfe}
\affiliation{Perimeter Institute for Theoretical Physics, Waterloo, Ontario, Canada, N2L 2Y5}

\author{Daniel Centeno D\'iaz}
\affiliation{\todo{affiliation}}

\author{Maria Ciudad Ala\~n\'on}
\affiliation{\todo{affiliation}}

\author{Pedro Nobrega Lauand}
\affiliation{\todo{affiliation}}

\author{Roberto Dobal Baldij\~ao}
\affiliation{\todo{affiliation}}

\author{Yujie Zhang}
\affiliation{\todo{affiliation}}

\date{\today}

\begin{abstract}
\todo{Abstract. One-paragraph summary: we quantify the secret-key rates achievable in prepare-and-measure contextuality-based key-distribution protocols against individual entangling-probe attacks with delayed, setting-dependent product measurements. We bound the eavesdropper's guessing probability both when she is constrained only by the operational theory (a linear program) and when she is constrained by quantum theory (a semidefinite program), and we provide an open-source pipeline that converts any such protocol specification into certified key rates together with dual-certificate witness inequalities.}
\end{abstract}

\maketitle

\section{Introduction}
\label{sec:intro}

% ------------------------------------------------------------------
% Paragraph 1: setting and motivation.
% ------------------------------------------------------------------
Prepare-and-measure protocols certify secrecy from the structure of the
observed statistics: Alice prepares one of a set of states labeled
$x$, Bob measures with a setting $y$ and obtains an outcome $b$, and the
pair distills a shared secret from the correlations $P(b|x,y)$ together
with the operational equivalences that the preparations and measurements
are promised to satisfy. When the statistics witness preparation
contextuality~\cite{Spekkens2005}, no eavesdropper whose side information
is mediated by a noncontextual (classical) model can predict the
outcomes perfectly, and a positive secret-key rate can survive.
\todo{Expand motivation; cite contextuality-based QKD and randomness
literature.}

% ------------------------------------------------------------------
% Paragraph 2: the security paradigm considered in this work.
% ------------------------------------------------------------------
In this work we analyze security against \emph{individual
entangling-probe attacks with delayed, setting-dependent product
measurements}. In each round, the eavesdropper Eve attaches an ancillary
probe to the system in transit and lets it interact, but she does not
measure immediately: she waits until the classical announcement of the
measurement setting $y$ (and, where relevant, of which rounds are
key-generating), and only then measures her probe, with a measurement
that may depend arbitrarily on $y$. Crucially, the attack is
\emph{individual}: Eve's probe interacts with each round separately and
is measured with a product (round-by-round) strategy, rather than a
collective or coherent attack across rounds. Within this paradigm we
consider two classes of adversary, matched to the two solver backends of
our accompanying code. The \emph{operational adversary} is constrained
only by the operational theory reproducing the observed statistics: her
side information is an arbitrary decomposition of the observed behavior
into subnormalized behaviors consistent with the promised operational
equivalences. Bounding this adversary's guessing probability is a linear
program (LP). The \emph{quantum adversary} is additionally constrained
by quantum theory: the joint state of message and probe, and Eve's
setting-dependent probe measurements, must admit a quantum realization.
Bounding this adversary is a noncommutative polynomial optimization
problem, which we relax to a semidefinite program (SDP) in the spirit of
the Navascu\'es--Pironio--Ac\'in hierarchy~\cite{NPA2008}, handling
nonprojective measurements through an explicit Naimark-unitary
construction. \todo{Sharpen the formal definitions and point to the
technical sections.}

% ------------------------------------------------------------------
% Paragraph 3: from protocols to key rates (the pipeline).
% ------------------------------------------------------------------
Our second contribution is a general-purpose, open-source pipeline that
converts a protocol specification into certified key rates. A protocol
is specified as a contextuality scenario: a behavior $P(b|x,y)$ (given
numerically, or generated from quantum states and effects, or from
states and effects of a generalized probabilistic theory) together with
its preparation and measurement operational equivalences, which the code
can also discover automatically. From this data the pipeline computes,
for each measurement setting $y$, Eve's optimal guessing probability
$\pguess(\cdot\,|y)$ for the raw key variable---held by Alice or by Bob,
as configured---under either adversary model, and converts these
guessing probabilities into uncertainties (min-entropy, or the tighter
reverse-Fano bound) and thence into asymptotic secret-key rates by
subtracting the error-correction cost between the honest parties.
Because every optimization is solved together with its dual, the
pipeline additionally returns human-readable \emph{witness
inequalities}: linear functionals of the observed behavior that certify
the reported bound and expose exactly which features of the statistics
power the security claim. \todo{Forward-reference the sections deriving
the LP, the SDP, and the key-rate formulas.}

% ------------------------------------------------------------------
% Paragraph 4: the where_key distinction and key-rate bookkeeping.
% ------------------------------------------------------------------
A structural distinction that our framework makes explicit is
\emph{which rounds generate key}. In the simplest protocols every
preparation setting $x$ and every measurement setting $y$ is used to
establish a key dit, and the key rate per experimental run coincides
with the key rate per key-generating round. In general, however, a
protocol may designate only a subset of preparations as key-generating
for each measurement setting---in our code, a nontrivial \wherekey{}
assignment listing, for each $y$, the values of $x$ whose rounds
contribute to the key---with the remaining rounds spent on parameter
estimation or discarded at sifting. This gives the notion of ``key rate
per experimental run'' a more nuanced meaning: it is the key rate per
key-generating round discounted by the fraction of rounds that are
key-generating at all, and optimizing the choice of \wherekey{} (which
our pipeline can do automatically) trades the quality of each key dit
against the rate at which key rounds occur. \todo{Give a small worked
example contrasting trivial and nontrivial \wherekey{}.}

% ------------------------------------------------------------------
% Paragraph 5: outline.
% ------------------------------------------------------------------
The remainder of the paper is organized as follows.
Section~\ref{sec:scenarios} fixes notation for contextuality scenarios
and operational equivalences. Section~\ref{sec:adversaries} formalizes
the two adversary models and derives the LP and SDP bounds on Eve's
guessing probability. Section~\ref{sec:keyrates} assembles the key-rate
formulas and the \wherekey{} bookkeeping. Section~\ref{sec:casestudies}
applies the pipeline to a collection of protocols, and
Section~\ref{sec:discussion} concludes.

\section{Contextuality scenarios and operational equivalences}
\label{sec:scenarios}

\todo{Notation: behaviors $P(b|x,y)$; preparation and measurement
operational equivalences; noncontextual models; quantum and GPT
realizations; automatic discovery of equivalences.}

\section{Adversary models}
\label{sec:adversaries}

\subsection{The operational adversary: a linear program}
\label{sec:lp}

\todo{Variables $P(b,e|x,y)$; data-consistency, operational-equivalence,
and positivity constraints; the guessing objective for a configurable
master-key holder (Alice or Bob); dual witness inequalities and their
$\ell_2$-refinement.}

\subsection{The quantum adversary: a semidefinite program}
\label{sec:sdp}

\todo{Entangling probe with delayed setting-dependent product
measurements; Naimark-unitary treatment of nonprojective measurements;
moment-matrix relaxation and NPA levels; U-only generator reduction;
complex-to-real embedding.}

\section{From protocols to key rates}
\label{sec:keyrates}

\todo{Guessing probability to min-entropy and reverse-Fano uncertainty;
error-correction cost between the honest parties; key rate per
key-generating round vs.\ per experimental run; the \wherekey{}
assignment and its automatic optimization.}

\section{Case studies}
\label{sec:casestudies}

\todo{BB84; the gbit; the $(3,2)$-PORAC; the 18-ray Cabello and 24-ray
Peres Kochen--Specker sets; hexagon and XZ-ring families; Bell-local
probability protocols. Tables of LP vs.\ SDP key rates.}

\section{Discussion}
\label{sec:discussion}

\todo{Beyond individual attacks; finite-size effects; outlook.}

\begin{acknowledgments}
\todo{Acknowledgments, including the Randomness Winter School 2026 and
funding agencies. This research was supported in part by Perimeter
Institute for Theoretical Physics.}
\end{acknowledgments}

\bibliography{references}

\end{document}
